\subsection{Lokförarnas behov och SFERA-protokollet}
Jämförelsen mellan lokförarnas informations- och kommunikationsbehov och SFERA-protokollet tyder på god förbättringspotential för SFERA som kan nås med enkla medel.
Framförallt beträffande kommunikation från lokförare till tågklarerare.
SFERA erbjuder dock ett fullgott kommunikationsprotokoll där en lokförares grundläggande kommunikationsbehov tillgodoses.
Trots detta finns förbättringsmöjligheter som med små förändringar i SFERA-protokollet kan leda till nytta och \lq{}quality of life\rq{}-förbättringar för lokförare.
Dessa förbättringar kan leda till att förarna får större förtroende för digitala verktyg vilket i sin tur kan leda till en brantare införandetakt av C-DAS hos järnvägsföretagen.
\paragraph{}
\hspace{-11pt}
Vikten av smidig kommunikation ska inte underskattas; digitala förarverktyg ska inte lösa problem som bäst avhjälps med andra medel.
Den största avvägningen i sammanhanget är huruvida datakommunikation faktiskt trumfar GSM-R-samtal.
Detta visade sig vara fallet för något av förbättringsförslagen vilket redovisas i~\ref{subsec:validering}.
Den bästa kommunikationsformen är kort, koncis och tydlig.
För lokförare saknas därför möjligheten att kommunicera i fritext både i SFERA och studiens förbättringsförslag; det skulle leda till en försämring av innevarande kommunikationsmedel (GSM-R-telefonsamtal).
Vidare har lokföraren många arbets- och distraktionsmoment som alla måste dela på förarens begränsade uppmärksamhetskvot.
Det finns emellertid många små korta samtal som äger rum idag vilka med enkelhet kan effektiviseras bort med datameddelanden i ett förbättrat SFERA-protokoll.
Att IT-artefakten banar väg för enkel och smidig kommunikation har därför varit avgörande.